\section{Related work}
\label{sec:related}

\paragraph{Channel pruning.}
Most pipelines start from one of three scores: the $\ell_1$ norm of a filter~\cite{li2017pruning}, the BatchNorm scale learned under a sparsity penalty~\cite{liu2017slimming}, or a first-order Taylor estimate of the loss change~\cite{molchanov2017pruning,molchanov2019importance}. He et al.~\cite{he2017channel} choose channels that best reconstruct the next layer's output and, for residual networks, prune only the inner convolutions of each block, which is the scheme we use. Discrimination-aware pruning~\cite{zhuang2018dcp} adds auxiliary classification losses to intermediate layers and keeps the channels that preserve them. DACIS asks a similar question with a cheaper statistic, the Fisher ratio of the pooled activation, and no extra training. These methods are evaluated on the classes they were trained for. Two earlier results suggest that such an evaluation tells us less than it seems to. Mittal et al.~\cite{mittal2018recovering} showed that removing $25$ to $50\%$ of the filters of VGG-16 and ResNet-50 at random and then fine-tuning matches ranked pruning, so that recovery, not the ranking, carries the result. Liu et al.~\cite{liu2019rethinking} found that a fine-tuned pruned network is no better than the same architecture trained from scratch. MetaPruning~\cite{liu2019metapruning} also says ``meta'', but there a hypernetwork predicts weights for candidate architectures; our meta stage is ordinary episodic training.

\paragraph{Is the ranking needed?}
Random choice has been strong in other pruning settings. Su et al.~\cite{su2020sanity} and Li et al.~\cite{li2022revisiting} report that randomly pruned networks with suitable layer-wise widths are competitive with criterion-based pruning, and Blalock et al.~\cite{blalock2020state} document how often pruning papers compare against weak baselines or inconsistently with one another. All of these evaluate on the classes the network was trained for. Tian et al.~\cite{tian2020mwnp} prune the network inside a meta-learner to reduce overfitting to the training tasks, but compare no channel scores. We did not find a comparison of channel scores on classes held out from the ones used to choose the channels. This paper makes one, with paired test episodes over several class partitions.

\paragraph{Few-shot classification.}
Prototypical networks~\cite{snell2017proto} label a query by the nearest class mean in embedding space, and MAML~\cite{finn2017maml} learns an initialisation that adapts in a few gradient steps. A network pre-trained with plain cross-entropy on the base classes is a strong baseline~\cite{chen2019closer,tian2020rethinking}, and episodic training on top of it with a cosine nearest-centroid head can add a little more~\cite{chen2021metabaseline}. Our stages S0 and M follow that recipe. The cross-domain benchmark of Guo et al.~\cite{guo2020bscd} has crop-disease and EuroSAT targets and reports that meta-learning methods often lose to simple fine-tuning when the target is far from the source. Fine-tuning a strong pretrained backbone is a strong baseline for few-shot classification~\cite{dhillon2020baseline,hu2022pushing}; those studies do not prune the backbone. Section~\ref{sec:results} asks whether the episodic stage still helps once the network is pruned.

\paragraph{Earth observation.}
Few-shot scene classification in remote sensing has been approached with metric learning on aerial benchmarks~\cite{li2021rsmetanet} and with meta-learning for land-cover classification~\cite{russwurm2020meta}. NWPU-RESISC45~\cite{cheng2017resisc} and the Sentinel-2~\cite{drusch2012sentinel2} patches of EuroSAT~\cite{helber2019eurosat} are standard scene and land-use benchmarks, which we split by class instead of by image. Pretraining on satellite imagery, for example with masked autoencoders~\cite{cong2022satmae}, is the natural alternative to an ImageNet initialisation; we did not test it. CloudScout~\cite{giuffrida2020cloudscout} runs a compact CNN on a Myriad~2 accelerator in the $\Phi$-Sat-1 mission~\cite{giuffrida2022phisat}. We do not target that hardware.

\paragraph{Plants.}
PlantVillage~\cite{hughes2015plantvillage} and PlantDoc~\cite{singh2020plantdoc} are common plant-disease datasets; Argüeso et al.~\cite{argueso2020fewshot} study few-shot recognition on PlantVillage with a $32/6$ class split, and we use a similar cross-class split.
